% !PRINTED 514
\weberpage{514}
%<<CONT
erty are divisible by one another, hence associated, and we accordingly set up the definition:

1. \gesperrt{The functional $\delta = \alpha x + \beta y + \cdots$ and every functional associated with it is called a greatest common divisor of $\alpha, \beta, \dots$}

If $\delta$ is a unit, we also say that $\alpha, \beta, \dots$ are without common divisor, for then there is, apart from the units, no integral functional which divides all of $\alpha, \beta, \dots$

2. \gesperrt{Two functionals $\alpha, \beta$ which have no common divisor, for which therefore $\alpha x + \beta y$ is a unit, are called relative prime functionals or coprime.}

From these definitions the following theorems result very simply:

3. \gesperrt{If integral functionals $\xi, \eta, \dots$ exist in $\Omega$ such that
\[ \alpha \xi + \beta \eta + \cdots = \varepsilon \]
is a unit, then the integral functionals $\alpha, \beta, \dots$ are without common divisor.}

For if $\alpha, \beta, \dots$ have a common divisor $\delta$, then (by \S.\,138) $\delta$ is also a divisor of $\alpha \xi + \beta \eta + \cdots$, and $\delta$ must therefore also be a unit.

4. \gesperrt{If $\alpha, \beta, \gamma$ are three integral functionals and $\alpha$ is relatively prime to $\beta$ and to $\gamma$, then $\alpha$ is also relatively prime to $\beta \gamma$.}

For by assumption, if $x, y, u, v$ are four variables,
\[ \varepsilon = \alpha x + \beta y, \quad \varepsilon_1 = \alpha u + \gamma v \]
are units. Accordingly
\[ \alpha (\alpha u x + \gamma v x + \beta u y) + \beta \gamma v y = \varepsilon \varepsilon_1 \]
is also a unit. But since $\alpha u x + \gamma v x + \beta u y$ and $v y$ are integral, it follows from Theorem 3. that $\alpha$ and $\beta \gamma$ are relatively prime.

To this is joined the proof of the following very important theorem:

5. \gesperrt{If $\alpha, \beta, \mu$ are integral functionals, $\alpha$ relatively prime to $\beta$ and $\alpha \mu$ divisible by $\beta$, then $\mu$ is divisible by $\beta$.}
